The \verb|intensity| function calculates the outgoing intensity of a single ray travelling through a system, represented by a \verb|Tree| object. Refraction and reflection flags
exist to give users the option to decide whether each phenomenon is considered. With both flags disabled, the ray streamlines through the system without any change in direction,
so it is a purely attenuation problem. With both flags enabled, the simulation becomes a binary recursion problem, as the ray splits into two rays every time it enters the interface
between two media. The calculated value is the sum of all rays that manage to escape the system.

Each of the test cases below compares the calculated intensity upon existing a sphere of radius $R = 2$, attenuation coefficient $\mu = 0.3$, and index of refraction $n = 1.3$ to the
corresponding known analytical solution.

\section{Streamline at normal incidence}
\begin{figure}[h!]
    \centering
    \includegraphics[scale=0.5]{Figures/attenuation_case1.png}
    \caption{Schematic plot of Test Case 1}
    \label{fig:att_case1}
\end{figure}

In a sphere, normal incidence requires that the ray enters along the sphere radius and crosses the center. This is the simpliest case, as the path length of the ray is exactly equal
to the diameter. Therefore, the intensity has been reduced to a factor of $A = e^{-2\mu r} = 0.301194$. The calculated value is exactly the same.

\section{Refraction and reflection at normal incidence}
\begin{figure}
    \centering
    \includegraphics[scale=0.4]{Figures/attenuation_case2.png}
    \caption{Schematic plot of Test Case 2.}
    \label{fig:att_case2}
\end{figure}

If a ray has normal incidence, the transmitted ray is not refraction, and the reflected direction is the opposite of incident. The reflectance and transmittance, defined as the ratio
between the reflected and transmitted ray intensity, respectively, to the incident intensity, is defined as:
\begin{align*}
    R &= \left( \frac{n_1 - n_2}{n_1 + n_2} \right)^2 \\
    T &= 1 - R
\end{align*}

In this case, the reflectance does not depend on whether the ray travels from medium 1 to 2 or vice versa. Taking $n_1 = 1.0$ to be vacuum and $n_2 = 1.3$, $R = 0.0170$ and
$T = 0.9830$.

Refer to Figure \ref{fig:att_case2} for a schematic plot. The events in this case are as followed:
\begin{enumerate}
    \item A normally incident ray enters the sphere with unit intensity.
    \item The ray splits into a reflected ray with intensity $R$ and a transmitted ray with intensity $T$. The reflected ray is collected (hence marked red), while the transmitted
    ray travels across the sphere and is attenuated to an intensity of $TA$.
    \item The uncollected ray splits into a reflected ray with intensity $TRA$ and a transmitted ray with intensity $T^2A$. The transmitted ray is collected, while the reflected ray
    travels across the sphere in the other direction and is attenuated to an intensity of $TRA^2$.
    \item The uncollected ray splits into a reflected ray with intensity $TR^2A^2$ and a transmitted ray with intensity $T^2RA^2$. The transmitted ray is collected, while the
    reflected ray travels across the sphere in the other direction and is attenuated to an intensity of $TR^2A^3$.
    \item Follow similar paths for Steps 5 and 6.
\end{enumerate}

The total outgoing intensity can be written as
\begin{align*}
    I &= R + T^2A + T^2RA^2 + T^2R^2A^3 + T^2R^3A^4 + \dots \\
    &= R + T^2A \sum_{i=0}^{\infty} (RA)^i \\
    &= R + \frac{T^2A}{1 - RA}
\end{align*}

Substituting in the values of $A$, $R$, and $T$, the true outgoing intensity is about $0.309545$. The calculated value is accurate to 6 significant digits and have a relative error
of $1.2716 \times 10^{-7}$.

\section{Streamline at non-normal incidence}
\begin{figure}[h!]
    \centering
    \includegraphics[scale=0.5]{Figures/attenuation_case3.png}
    \caption{Schematic plot of Test Case 3}
    \label{fig:att_case3}
\end{figure}
In this test case, the chosen direction is $\hat{\imath}$, incident upon the sphere at position $\displaystyle \left[\frac{\sqrt{3}r}{2}, \frac{r}{2}, 0 \right]^T$. The path length
is therefore $\sqrt{3}r$, and the intensity is attenuated to $e^{-\sqrt{3}\mu r} = 0.353727$ upon exiting the sphere. Again, the value calculated by the \verb|intensity| function is
identical to the analytical solution.